\documentclass[11pt,a4paper]{article}
\usepackage{brainteaser}

\begin{document}
\brainteasertitle{16}{Motifs dans les décimales}

\begin{statementbox}
On choisit uniformément un nombre $X$ dans $[0,1]$ et on lit successivement
les chiffres de son développement décimal:
\[
X=0,D_1D_2D_3\ldots
\]

\begin{enumerate}
  \item Combien de décimales faut-il lire en moyenne avant de voir apparaître
  le motif $222$? Généraliser au motif formé de $n$ chiffres $2$
  consécutifs.
  \item Combien de décimales faut-il lire en moyenne avant de voir trois
  chiffres identiques consécutifs, sans imposer leur valeur? Généraliser à
  une suite de $k$ chiffres identiques quelconques.
\end{enumerate}

On compte toutes les décimales lues, y compris celles qui terminent le motif.
\end{statementbox}

\begin{hintsbox}
\begin{enumerate}
  \item Les chiffres $D_i$ sont indépendants et uniformes sur
  $\{0,1,\ldots,9\}$.
  \item Pour chercher $n$ chiffres $2$, définir un état $j$ égal au nombre de
  chiffres $2$ consécutifs observés à la fin de la suite courante.
  \item Après un chiffre différent de $2$, cet état revient à $0$.
  \item Pour chercher $k$ chiffres identiques quelconques, mémoriser seulement
  la longueur de la série courante.
  \item Dans ce second problème, un chiffre différent du précédent ne ramène
  pas à un état vide: il commence immédiatement une nouvelle série de
  longueur $1$.
\end{enumerate}
\end{hintsbox}

\correction

À un ensemble de probabilité nulle près, un nombre uniforme sur $[0,1]$
possède un développement décimal unique. Ses chiffres $D_1,D_2,\ldots$ sont
indépendants et chacun prend les dix valeurs $0,1,\ldots,9$ avec probabilité
$1/10$. Le choix entre les deux écritures des nombres décimaux, par exemple
$0{,}5000\ldots$ et $0{,}4999\ldots$, ne change donc pas les espérances.

\paragraph{Le motif $222$.}
Il faut tenir compte des recouvrements possibles du motif. Notons $E_j$ le
nombre moyen de chiffres supplémentaires à lire lorsque les derniers chiffres
observés sont exactement $j$ chiffres $2$ consécutifs, avec $0\leq j\leq3$.
L'état $3$ est absorbant, donc $E_3=0$.

Depuis l'état $0$, le prochain chiffre est un $2$ avec probabilité $1/10$;
sinon on reste dans l'état $0$. Ainsi,
\[
E_0=1+\frac1{10}E_1+\frac9{10}E_0.
\]
Depuis les états $1$ et $2$, un nouveau $2$ fait progresser d'un état, tandis
que tout autre chiffre ramène à l'état $0$:
\[
E_j=1+\frac1{10}E_{j+1}+\frac9{10}E_0,
\qquad j=1,2.
\]
La première équation donne $E_0=10+E_1$. En utilisant ensuite les deux autres
équations et $E_3=0$, on obtient
\[
E_0=10+100+1000=1110.
\]

\begin{resultbox}
\[
\boxed{\E[T_{222}]=1110.}
\]
Il faut lire en moyenne $1110$ décimales avant la première apparition de
$222$.
\end{resultbox}

\paragraph{$n$ chiffres $2$ consécutifs.}
Conservons les états $0,1,\ldots,n$, avec $E_n=0$. Posons
\[
A_j=E_0-E_j.
\]
Les mêmes équations donnent $A_0=0$ et
\[
A_{j+1}=10(A_j+1),
\qquad 0\leq j<n.
\]
Par récurrence,
\[
A_j=10+10^2+\cdots+10^j.
\]
Or $A_n=E_0-E_n=E_0$. Par conséquent,
\[
\boxed{
\E[T_{2,n}]
=10+10^2+\cdots+10^n
=\frac{10(10^n-1)}9.}
\]
Pour $n=1,2,3,4$, les espérances sont respectivement
$10$, $110$, $1110$ et $11110$.

\paragraph{Trois chiffres identiques quelconques.}
Cette fois, la valeur du chiffre répété n'est pas fixée. Après avoir lu le
premier chiffre, une série de longueur $1$ est déjà commencée. Notons $F_j$ le
nombre moyen de chiffres supplémentaires nécessaires lorsque la série finale
courante a longueur $j$, avec $1\leq j\leq3$. On a $F_3=0$.

Pour $j=1,2$, le chiffre suivant est identique au précédent avec probabilité
$1/10$, ce qui prolonge la série. Avec probabilité $9/10$, il est différent:
la série recherchée échoue, mais ce nouveau chiffre commence aussitôt une
nouvelle série de longueur $1$. Ainsi,
\[
F_j
=1+\frac1{10}F_{j+1}+\frac9{10}F_1,
\qquad j=1,2.
\]
Ces équations donnent
\[
F_1=10+100=110.
\]
Il faut ajouter le tout premier chiffre, lu avant d'entrer dans l'état $1$.
Le nombre moyen total de décimales est donc
\[
\boxed{1+F_1=111.}
\]

\paragraph{$k$ chiffres identiques quelconques.}
Avec les états $1,2,\ldots,k$, posons
\[
B_j=F_1-F_j.
\]
On a $B_1=0$ et les équations précédentes donnent
\[
B_{j+1}=10(B_j+1),
\qquad 1\leq j<k.
\]
Donc
\[
F_1=B_k=10+10^2+\cdots+10^{k-1},
\]
car $F_k=0$. En ajoutant le premier chiffre, on trouve
\[
\boxed{
\E[T_k]
=1+10+10^2+\cdots+10^{k-1}
=\frac{10^k-1}{9}.}
\]
Pour $k=1,2,3,4$, cela donne $1$, $11$, $111$ et $1111$ décimales en
moyenne.

\paragraph{Comparaison.}
Chercher $k$ occurrences consécutives d'un chiffre fixé demande en moyenne
exactement dix fois plus de chiffres que chercher $k$ chiffres identiques de
valeur quelconque:
\[
\frac{10(10^k-1)}9
=10\,\frac{10^k-1}{9}.
\]

\paragraph{Dans une base quelconque.}
En base $b$, les mêmes raisonnements remplacent $10$ par $b$. Pour un chiffre
fixé répété $n$ fois et pour une valeur quelconque répétée $k$ fois, on obtient
respectivement
\[
\boxed{
\E[T_{a,n}]=\frac{b(b^n-1)}{b-1},
\qquad
\E[T_k]=\frac{b^k-1}{b-1}.}
\]

\paragraph{Pourquoi pas $1000$ et $100$?}
Une fenêtre de trois chiffres vaut $222$ avec probabilité $1/1000$, et elle
est formée de trois chiffres identiques quelconques avec probabilité $1/100$.
On ne peut toutefois pas inverser directement ces probabilités: deux fenêtres
successives partagent deux chiffres et ne sont donc pas indépendantes. Les
états précédents mémorisent précisément l'information conservée d'une fenêtre
à la suivante.

\end{document}
